% Level 2, Week 12 -- Showcase and Celebration
\lesson{12}{Showcase and Celebration}{Share what you built, explain how it works, and see how far you have come.}
\lessoninfo{Concepts: explaining design (R), feedback (R), all Level 2 (R)}{Tools: Thonny, feedback cards, Assessments booklet, certificates}{CSTA: 2-AP-19, 2-AP-15, 2-IC-23}

\begin{objectives}
  \item The student can present the team's project and explain one design decision and one piece of code (a function, the loop, the list, the dictionary or the file). (L2.G8; L2.W12.1)
  \item The student can give kind, specific feedback to other makers during a gallery walk. (L2.G8; L2.W12.2)
  \item The student can show what they learned in Level 2 on the exit assessment: files, tests, events and functions. (L2.G3, L2.G5, L2.G6, L2.G7; L2.W12.3)
  \item The student can explain how to share a project safely: family only, nicknames, no personal information on screen. (L2.G9; L2.W12.4)
\end{objectives}

\begin{materials}
  \item Print: the Exit assessment in the Assessments booklet (one per student) and its answer page from the answer key (for you only); workbook chapter 12; homework 12; one Level 2 certificate per student (the Certificate document in \texttt{extras}); 4 feedback cards per student.
  \item On each Mac: the team's project open in Thonny, with a backup copy saved (\menu{File > Save as…} \texttt{mygame-showcase.py}) in case of last-minute accidents.
  \item Optional: invite family members to the showcase (minutes 28--50). Snacks make it a celebration!
  \item Pencils, a timer. Prep: 10 minutes, mostly printing.
\end{materials}

\needspace{16\baselineskip}
\section{At a glance}
\begin{timeplan}
  \timeslot{0--8}{Warm-up}{Freeze and rehearse: last save, fill in the presentation card, set up each computer as a gallery station.}
  \timeslot{8--28}{Concept}{Concept check: the Exit assessment in the Assessments booklet (20 minutes, individual, quiet).}
  \timeslot{28--38}{Guided demo}{Showcase presentations: you model a 1-minute pitch with Star Catcher, then each team presents (3 minutes + 1 question).}
  \timeslot{38--50}{Hands-on build}{Gallery walk: everyone plays every project and leaves feedback cards. Last minutes: one live tweak or free play.}
  \timeslot{50--60}{Wrap-up}{Read feedback cards, certificates, reflection, Level 3 preview, homework 12.}
\end{timeplan}

\begin{teachernote}[Why the phases look different today]
The \enquote{Concept} slot is a concept check (the exit assessment) while minds are fresh; the
celebration comes afterwards. Unfinished projects are normal: teams present what works and say what
they would add next.
\end{teachernote}

\section{Warm-up (0--8 min)}
\begin{unplugged}[Freeze and rehearse]
Say \enquote{Code freeze!}: no more new features (only fixing a crash is allowed). Each student fills in
the presentation card in workbook chapter 12 (What is it? How does one part work? What would you add
next?) and practices the pitch once to their partner, quietly, in one minute. Then set up the gallery:
each computer shows its project, with a stack of feedback cards next to it.
\end{unplugged}

\section{Concept (8--28 min)}
\begin{teacherscript}
“Before the party, one last challenge. This is not a test you can fail: it shows you, and me, how much you
learned since week 1. Read each question carefully. If you get stuck, skip it and come back. It is
pencil and paper only: no computer, no notes and no partner this time. You have 20 minutes.”
\end{teacherscript}
Hand out the Exit assessment from the Assessments booklet. Keep the room quiet and the laptops closed.
At 20 minutes, collect papers, even unfinished ones.
Score them later with the answer key (assessment section); do not score during the session.

\begin{analogy}[A video game save screen]
At the end of a level, a game shows your stats: coins collected, time, secrets found. The exit
assessment is that stats screen for Level 2. It is not the boss fight. It shows what you have
unlocked so far.
\end{analogy}

\section{Guided demo (28--38 min)}
\begin{workedexample}[Model a 1-minute pitch]
Go first, so students see that a pitch is short and simple. Run Star Catcher (or any team's game, with
permission) and say:
\begin{enumerate}
  \item \textbf{What it is:} \enquote{This is Star Catcher. You catch falling stars with a basket and dodge meteors.}
  \item \textbf{How one part works:} \enquote{A timer runs the game loop many times a second. Each time, it moves every star in a list, and a logic function returns True when a star is close enough to the basket. That function has tests.}
  \item \textbf{A design decision:} \enquote{Stars and meteors have different shapes, not only different colors, so everyone can tell them apart.}
  \item \textbf{What next:} \enquote{I would add a slow mode for beginners.}
\end{enumerate}
Then each team presents for 3 minutes and takes one question. Audience questions to suggest:
\enquote{Which part was hardest?} \enquote{Show us one of your tests.} \enquote{What does your file save?}
\end{workedexample}

\begin{pairroles}[Presenter roles]
\textbf{2 per team:} one is the \textbf{demo driver} (plays the game, clicks), the other is the
\textbf{explainer}; swap halfway so both speak. \textbf{1 student:} the student does both; you ask the
question. \textbf{3:} driver, explainer, and \textbf{code guide} who shows one function on screen.
\textbf{4:} two teams, each presents to the other (and to any guests).
\end{pairroles}

\section{Hands-on build (38--50 min)}
\begin{safety}[Showcase safety]
Before guests play, check each screen and file together:
\begin{checklist}
  \item High-score files and saved designs use nicknames only (open \texttt{highscore.txt} and look).
  \item No real names, school names, addresses, ages or photos appear in the code, the window or the title.
  \item The showcase is for family. If someone wants to take a photo or video, film the screen, not faces, and ask everyone first.
  \item Nothing gets uploaded or posted anywhere without a parent deciding together with the student.
\end{checklist}
Say: \enquote{You can be proud of your work and still keep yourself private. Your nickname is your artist name.}
\end{safety}

\textbf{Gallery walk (10 minutes).} Everyone moves from station to station, plays each project for
2--3 minutes and fills in a feedback card: \textbf{two stars} (two specific things that are good) and
\textbf{a wish} (one kind idea). Leave the card face down next to the computer. Makers do not explain or
defend during the walk; they just watch people play (that is great playtest data!). Guests can fill in
cards too.

\begin{freebuild}[Last two minutes: live tweak]
Each team may make one tiny live change suggested by a card (a color, a speed, a message) while
others watch, or simply play each other's games.
\end{freebuild}

\section{Wrap-up (50--60 min)}
Makers read their feedback cards (2 minutes). Then hand out certificates one by one: say one specific
thing each student did well this level (\enquote{You kept going when your collision test failed, and
you found the bug}). Read the Level 3 preview: \enquote{In Level 3 you will save your code with
\textbf{git save points}, build your own \textbf{classes} (blueprints for objects), and the big example
is \textbf{Maze Muncher}, a maze game where a ghost finds you with a clever search.} Hand out homework 12:
a reflection and the holiday \enquote{keep coding} menu.

\begin{formativecheck}[What to look for]
\textbf{In the pitch:} the student names a real function or the loop and says what goes in and what
comes out; mentions at least one data structure (list, dictionary or file).
\textbf{Exit-ticket question:} \enquote{Which Level 2 tool would you use to make a program remember
something after it closes?}
\textbf{Answer:} a file (\texttt{open} with \texttt{"w"} or \texttt{"a"} to save, \texttt{"r"} to load).
\end{formativecheck}

\begin{reflection}
What are you proudest of in your project? What did you learn about working in a team? Which
Level 2 skill do you want to use again, and what do you want to build next?
\end{reflection}

\section{Notes for the teacher}
\begin{commonmistakes}
  \item \textbf{The project crashes during the demo.} Open the backup copy and carry on; say \enquote{that is live coding!} Every professional demo has had this happen.
  \item \textbf{Apologizing for unfinished parts} (\enquote{it is bad, it does not work}). Coach: \enquote{Say what works, then what you would add next.}
  \item \textbf{Vague feedback} (\enquote{good}, \enquote{cool}). Ask: \enquote{Which part exactly? What did you do when you smiled?}
  \item \textbf{Comparing projects} (\enquote{ours is better}). Remind the group: every team started from a different idea; we compare with our own week 1 selves, not with each other.
\end{commonmistakes}

\begin{differentiation}
\supportpath Give the presentation card with sentence starters already written (\enquote{My project is…},
\enquote{This function…}); let a shy student present to you alone, or let the partner drive while they
point. For the assessment, read questions aloud if reading slows them down (note it on the paper).
\extensionpath Ask the student to show and explain one of their tests and the bug it caught, or to
write a short \enquote{How to play} text file (\texttt{README.txt}) for players.
\end{differentiation}

\begin{stretchchallenge}[Version 2 plan]
Write a \enquote{version 2} plan for the project: three new features in pseudocode, which list or
dictionary each would need, and one new test for each. Keep it for Level 3 week 1.
\end{stretchchallenge}

\begin{teachernote}[If you are short on time]
Keep the exit assessment and the certificates. Shorten the presentations to 2 minutes each and do the
gallery walk with one card per student. The reflection can move to homework 12.
\end{teachernote}
